\begin{statement}{C4f}{Proved}{An unrestricted geometric replacement for CDP20 Theorem 2.2}
Let $f:Y\to\PP^1$ satisfy the integral CDP hypotheses.
For each critical full fibre with
irreducible nonnormal reduced support $S$, let
$\eta_S:T_S\to S$ be its normalization. Then a general
line bundle $L\in\Pic(Y)$ satisfies
\[
 H^q(Y,T_Y\otimes L)=0\quad(q\ne1)
 \tag{C.VC.1}
\]
exactly when, for every such $S$, the holomorphic pullback
homomorphism
\[
 \Pic(Y)\longrightarrow\Pic(T_S)
 \tag{C.VC.2}
\]
is nonzero. Equivalently, its differential
$H^1(Y,\OO_Y)\to H^1(T_S,\OO_{T_S})$ is nonzero.
In this case $\chi(Y,T_Y\otimes L)\le0$ and $c_3(Y)\le0$.
If one map (C.VC.2) is zero, then
$H^2(Y,T_Y\otimes L)\ne0$ for every $L\in\Pic(Y)$.
All fibre multiplicities and all reduced-support
singularities are permitted in this criterion.
\end{statement}
\begin{proof}
The classification and exclusions cited in C4e leave
exactly the two cohomological cases used in that proof:
torus general fibre, or $h^1(\OO)=1$, $h^2(\OO)=0$.
The Picard group of $Y$ is the connected exponential image
of $H^1(Y,\OO_Y)$. A holomorphic group homomorphism out of
this group is zero precisely when its differential is zero.
For each fixed critical full fibre, the ambient-section
nonvanishing locus is closed analytic in this group, as
in C4c. A.3.4 makes it proper when the support is smooth,
and A.3.5 makes it proper when the support is reducible.
A connected normal support is irreducible and is smooth
by A.5.2, so these cases leave just the supports in the
statement.

Suppose the nonvanishing locus for $F=mS$ is the whole
Picard group. The proof of C4e, beginning with its
fixed-fibre persistence assertion, proves that (C.VC.2)
is zero. Its steps use the actual component filtration,
the inverse-Jacobian kernel, the Gauduchon bound, the
effective conductor comparison, the coherent conductor
genus bound, and the explicit constant-transverse-type
calculations. Each is independent of the Euler number
once persistence has been given.

Conversely, suppose (C.VC.2) is zero. The last nonzero
ideal layer of the full fibre is an actual subsheaf
\[
 (\Omega_Y^1\otimes M\otimes\OO_Y(-(m-1)S))|_S
 \hookrightarrow
          (\Omega_Y^1\otimes M)|_{mS}.
 \tag{C.VC.3}
\]
The canonical normalization submodule in A.2.2 gives
an injection into the left side from
\[
 \eta_{S*}\bigl(\eta_S^*((M\otimes\OO_Y(-mS))|_S)\bigr).
 \tag{C.VC.4}
\]
Here the first conormal factor $\OO_Y(-S)|_S$ combines
with the displayed $(m-1)$st ideal layer; this comparison
retains every nilpotent layer needed for the injection.
Since $mS$ is a full fibre, $\OO_Y(-mS)|_S$ is a trivial
line, with its one-dimensional scalar factor determined
by a local base parameter. The pullback of $M$ to $T_S$
is trivial by (C.VC.2). Hence (C.VC.4) has a nonzero global
section for every $M$, and its injections prove that the
full-fibre nonvanishing locus is the entire Picard group.
This also proves the converse directly for arbitrary $m$.

We have proved an exact geometric characterization of
every improper critical-fibre locus. For torus general
fibre, applying C4c gives the equivalence with generic
degree-two vanishing. In the second case A.5.13 identifies
$\Pic(Y)$ with the transgression line, so properness of
all these loci is exactly the premise of C4d and gives
generic degree-two vanishing there as well.
The endpoint vanishings of C.1 then give (C.VC.1) and the
index inequalities.
For the last assertion in the torus case, generic non-torsion on every
smooth torus fibre makes $R^2f_*(T_Y\otimes L)$ a
torsion sheaf for a dense set of twists, by D.1.11.
The fixed full-fibre obstruction makes it nonzero.
Its nonzero global sections are a quotient of
$H^2(Y,T_Y\otimes L)$ by the Leray sequence.
Thus degree-two cohomology is nonzero on a dense set.
Its proper semicontinuity locus is closed analytic in
the connected Picard family, so it is nonzero for every
twist. In the second case use instead the actual twists
$L_r=K_Y\otimes M_r^{-1}$ for real irrational $r$.
A.4.9 and the smooth-fibre conormal sequence make their
$R^2f_*(T_Y\otimes L_r)$ torsion. The fixed full-fibre
obstruction again makes that torsion sheaf nonzero,
so Leray gives $H^2(Y,T_Y\otimes L_r)\ne0$.
These twists have accumulation points in the complex
line $\Pic(Y)=\C\beta$. A closed complex analytic subset
of that line containing them is the whole line.
This proves all-twist nonvanishing in the second case
and completes both implications of the criterion.
\end{proof}


\subsection{C.5. The normalization quotient and every tangent twist}

% C5a from part_C2_C4_C5.tex

